\documentclass[11pt,a4paper]{article}
\usepackage[utf8]{inputenc}
\usepackage[T1]{fontenc}
\usepackage[spanish]{babel}
\usepackage[margin=2.5cm]{geometry}
\usepackage{array}
\title{GA-TL-02: Revisión de diseño de CLAB}
\author{}
\date{}
\begin{document}
\maketitle

\section*{Estado del documento}
Esqueleto de preparación. La revisión con el docente y la integración de la tabla de servicios están pendientes. Este archivo no es la entrega final.

\section*{Preparación de la revisión}
Las respuestas se elaborarán a partir del diagrama y las decisiones acordadas por el equipo.
\begin{enumerate}
\item ¿Qué ocurre si cae un servicio durante una operación y cómo se recupera?
\item ¿Qué responsabilidad y datos pertenecen a cada servicio? ¿Por qué mantener separados dos servicios relacionados?
\item ¿Cómo se mantiene la consistencia entre bases de datos y se evitan operaciones duplicadas?
\item ¿Qué canales son sincrónicos o asincrónicos y qué requisito justifica la tecnología elegida?
\item ¿Qué sucede si falla el API Gateway o el mecanismo de mensajería propuesto?
\end{enumerate}

\section*{Tabla de retroalimentación}
Sin observaciones registradas: la revisión con el docente está pendiente. Las filas se completarán únicamente con lo ocurrido en esa revisión.

\begin{center}
\renewcommand{\arraystretch}{1.4}
\begin{tabular}{|p{0.29\linewidth}|p{0.29\linewidth}|p{0.29\linewidth}|}
\hline
\textbf{Observación} & \textbf{Cambio realizado en el diagrama} & \textbf{Razón si no se aplicó} \\
\hline
\rule{0pt}{2cm} & & \\
\hline
\end{tabular}
\end{center}

\section*{Asistencia utilizada}
Estructura inicial y preguntas de preparación elaboradas con asistencia de OpenAI Codex. No se ha realizado ni simulado la revisión docente.
\end{document}
